\uuid{Lv2q}
\titre{Série télescopique et décomposition en éléments simples}

\niveau{L1}
\module{Séries numériques}
\chapitre{Séries à termes positifs}
\sousChapitre{Séries télescopiques}
\theme{décomposition en éléments simples, série télescopique, calcul de somme}
\auteur{}
\datecreate{2026-05-11}
\organisation{AMSCC}
\difficulte{2}

\contenu{

\texte{
On considère la série
$$
\sum_{n\geq 1}\frac{1}{(2n-1)(2n+1)}.
$$
}

\begin{enumerate}

\item \question{Montrer que, pour tout $n\geq 1$,
$$
\frac{1}{(2n-1)(2n+1)}=\frac{1}{2}\left(\frac{1}{2n-1}-\frac{1}{2n+1}\right).
$$
}
\reponse{
On effectue une décomposition en éléments simples. On cherche $A, B \in \mathbb{R}$ tels que
$$
\frac{1}{(2n-1)(2n+1)}=\frac{A}{2n-1}+\frac{B}{2n+1}.
$$
En réduisant au même dénominateur :
$$
\frac{A}{2n-1}+\frac{B}{2n+1} = \frac{A(2n+1)+B(2n-1)}{(2n-1)(2n+1)}.
$$
On identifie :
$$
A(2n+1)+B(2n-1)=1 \quad \text{pour tout } n\geq 1.
$$
En posant $n=\tfrac12$ : $2A=1$ donc $A=\tfrac12$.
En posant $n=-\tfrac12$ : $-2B=1$ donc $B=-\tfrac12$.

On obtient bien
$$
\frac{1}{(2n-1)(2n+1)}=\frac{1}{2}\left(\frac{1}{2n-1}-\frac{1}{2n+1}\right).
$$
}

\item \question{En déduire que la série $\displaystyle\sum_{n\geq 1}\frac{1}{(2n-1)(2n+1)}$ converge, et calculer sa somme.}
\indication{Écrire la somme partielle d'ordre $N$ et faire tendre $N$ vers $+\infty$.}
\reponse{
D'après la question précédente, la somme partielle d'ordre $N$ vaut
$$
S_N = \sum_{n=1}^N \frac{1}{(2n-1)(2n+1)}
= \frac{1}{2}\sum_{n=1}^N\left(\frac{1}{2n-1}-\frac{1}{2n+1}\right).
$$
La série est télescopique : les termes consécutifs se simplifient deux à deux,
$$
S_N = \frac{1}{2}\left(1 - \frac{1}{2N+1}\right).
$$
Comme $\dfrac{1}{2N+1}\xrightarrow[N\to+\infty]{}0$, on obtient
$$
\sum_{n\geq 1}\frac{1}{(2n-1)(2n+1)}=\frac{1}{2}.
$$
}


\end{enumerate}

}
